\section[$\Bun_G$]{$\Bun_G$（12 月 11 日）}

\srcpage{66}

\subsection{遗留问题}

我们需要补完上次的一个论证。我们当时试图证明：

\begin{theorem}
映射 $|\Bun_G|\to B(G)$ 是双射且连续。
\end{theorem}

为完成证明，尚缺下述结论：

\begin{theorem}
设 $\Bun^1_G\subset\Bun_G$ 是所有满足如下条件的 $G$-主丛 $\cE/X_S$ 构成的子叠：在所有几何点\footnote{这里我们总是指 $C$ 完全且代数闭，且 $C^+\subset C$ 是赋值子环。} $\Spa(C,C^+)\to S$ 处，$\cE|_{X_{\Spa(C,C^+)}}$ 平凡。则 $\Bun^1_G\subseteq\Bun_G$ 是开子叠，且 $\Bun^1_G\cong[\mathrm{pt}/G(E)]$。
\end{theorem}

\begin{remark}
我们已在 §10.3.1 对 $\mathrm{GL}_n$ 陈述过这一点。注意这个开性由「$|\Bun_G|\to B(G)$ 是双射且连续」这一陈述立即推出，因为平凡等晶体显然在 $B(G)$ 中定义一个开点。
\end{remark}

策略是通过 Tannaka 形式化（Tannakian formalism）化归到 $\mathrm{GL}_n$ 的情形。但这并非形式推理。

\begin{proof}
我们知道 Newton 映射 $\nu$ 半连续，而 $\Bun^1_G$ 包含在 $\nu=0$ 的轨迹中，后者是开的。因此可以把注意力限制在使 $S$ 的每个几何点都有 $\nu\equiv 0$ 的 $G$-主丛 $\cE$ 上。

在此情形，对一切表示 $\rho\colon G\to\mathrm{GL}(V)$，可考虑 $\rho_*\cE\in\mathrm{VB}(X_S)$。它半稳定、斜率为 $0$。因此它等价于 $S$ 上的（pro-étale）$\ul{E}$-局部系统（见前文关于 pro-étale 局部系统的推论）。

这仍然复杂，所以我们作 pro-étale 局部化，化归到 $S$ 严格完全间断（strictly totally disconnected）的情形。令 $A=\operatorname{Cont}(|S|,E)=\operatorname{Cont}(\pi_0(S),E)$。由于 $S$ 严格完全间断，$E$-局部系统的范畴通过整体截面函子等价于 $\Proj(A)$（有限投射 $A$-模的范畴）。

总之，$\cE$ 定义了正合 $\otimes$-函子
\[
\Rep_E(G)\to\{\mathrm{VB}(X_S)\text{ 处处半稳定、斜率为 }0\}\cong\Proj(A).
\]
这与 $\Spec A$ 上 $G$-主丛 $\cF$ 的数据相同，且 $G$-主丛 $\cE$ 的平凡性等价于 $G$-主丛 $\cF$ 的平凡性。所以只需看出：若 $\cF$ 在对应于 $s\in S$ 的点 $\Spec E\hookrightarrow\Spec A$ 处平凡（在 $s$ 处求值给出映射 $A\to E$），则它在拉回到
\[
\Spec\operatorname{Cont}(U,E)\subset\underbrace{\Spec\operatorname{Cont}(|S|,E)}_{A}
\]
之后平凡，其中 $U\ni s$ 是某个开（因而也闭）邻域。这由以下两个事实推出：

\begin{enumerate}[label=(\arabic*)]
\item 局部环
\[
\varinjlim_{U\ni s}\operatorname{Cont}(U,E)
\]
沿 $s$ 处求值的核是 hensel 的，剩余域为 $E$。
\item 若 $(B,I)$ 是 hensel 对，则 $H^1_{\et}(\Spec B,G)\hookrightarrow H^1_{\et}(\Spec B/I,G)$，即 $B$ 上任何在 $B/I$ 上分裂的 $G$-主丛已在 $B$ 上分裂。
\end{enumerate}
\end{proof}

\srcpage{67}

为结束证明，用 (2) 得到在 hensel 局部环上的分裂，再用 (1) 把它展布到邻域。关于 (1)，注意 $A$ 是一个 Tate 代数，其 adic 谱为 $\pi_0(S)$。该余极限于是是这个解析 adic 空间的局部环，而一个一般引理说解析 adic 空间的局部环总是 hensel 的。（这类似于复解析簇的情形，那里局部环由反函数定理是 hensel 的。）

\subsection{关于局部 Shimura 簇的插叙}

「局部 Shimura 簇」（参看 Rapoport–Viehmann）是 Shimura 簇的 $p$ 进（因而「局部」）类比，并通过单值化结果与 Shimura 簇相联系。例如，复 Shimura 簇是 Hermitian 对称域对算术群作用的商，而 Shimura 簇中的某些轨迹可以写成 $p$ 进对称空间对算术群作用的商。这一现象最早由 Cerednik 在 1970 年代研究。1980 年代，Rapoport–Zink 考虑了 PEL 型 Shimura 簇；在此情形，相应的「局部 Shimura 簇」是 $p$-可除群的模空间，推广了 Drinfeld 的一个想法。它们现在称为「Rapoport–Zink 空间」。

现在我们知道，一般的 Shimura 簇不承认模论式的描述。[RV14] 猜想这类对象的局部类比存在，而结果表明它们可以用我们已经建立的机器构造（参看 [Berk]）。

\subsubsection{局部 Shimura 数据}

我们给出定义局部 Shimura 簇所用的抽象群论数据。通常考虑 $E=\Qp$，但我们的形式化允许更一般的 $E$。

\begin{definition}
局部 Shimura 数据（local Shimura datum）是一个三元组 $(G,[b],\{\mu\})$，其中
\begin{itemize}
\item $G/E$ 是约化群；
\item $\{\mu\colon\GG_m\to G_{\breveE}\}$ 是极小余特征的共轭类；
\item $[b]\in B(G)$。
\end{itemize}
\end{definition}

为了使相应的局部 Shimura 簇非空，我们需要求
\[
[b]\in\underbrace{B(G,\mu)}_{\text{有限}}\subset B(G),
\]
其中 $B(G,\mu)$ 是由显式组合判据给出的有限子集。我们稍后给出该判据的几何表述。

\subsubsection{局部 Shimura 簇}

局部 Shimura 簇是一个塔
\[
(\mathcal M_{(G,b,\mu),K})_{K\ \text{紧开}\subseteq G(E)}
\]
由 $\breveE$ 上的光滑刚性解析簇组成（外加一个 Weil 下降数据，即与 $\breveE$ 上 Frobenius 拉回的同构），并带有相容的 étale 周期映射
\[
\pi_K\colon \mathcal M_{(G,b,\mu),K}\to\mathcal{F}\ell_\mu/\breveE.
\]
这里 $\mathcal{F}\ell_\mu$ 参数化 $G$ 的抛物子群，其共轭类由 $\mu$ 决定。$\pi_K$ 的非空几何纤维同构于 $G(E)/K$。

\begin{example}[Drinfeld 情形]
\footnote{讨论差一个符号。}若 $G=D^{\times}$，$D/E$ 是四元数代数，则可取 $\mu\colon\GG_m\to G_{\breveE}\cong\mathrm{GL}_2$ 为
\[
t\mapsto\begin{pmatrix}t&\\&1\end{pmatrix},
\]

\srcpage{68}

再取 $b$ 为斜率为 $1/2$ 的「基本」等晶体。（这是 $B(G,\mu)$ 的唯一元素。）于是 $\mathcal{F}\ell_\mu=\PP^1$（作为 adic 空间），而周期映射 $\pi_K$ 将取值于「Drinfeld 上半空间」$\Omega^2=\PP^1\setminus\PP^1(E)$。$\mathcal M_K$ 将是 Drinfeld 上半空间的 Drinfeld 覆盖。

对 $n\ge 2$ 的 $\mathrm{GL}_n$ 有类似的例子。取 $G$ 为不变量为 $1/n$ 的除法代数的单位群，取 $b$ 为斜率为 $1/n$ 的基本元素。目标将是 $\PP^{n-1}$ 去掉所有有理超平面。
\end{example}

\begin{example}[Lubin–Tate 情形]
\footnote{讨论差一个符号。}设 $G=\mathrm{GL}_n$，$\mu\colon\GG_m\to G$ 为
\[
t\mapsto\begin{pmatrix}
t & & & \\
& 1 & & \\
& & \ddots & \\
& & & 1
\end{pmatrix},
\]
并取 $b$ 为斜率为 $1/n$ 的基本元素。在此情形 $\mathcal{F}\ell_\mu=\PP^{n-1}/\breveE$，而周期映射
\[
(\mathcal M_K)_{K\subseteq\mathrm{GL}_n(E)}\to\mathcal{F}\ell_\mu
\]
就是 Gross–Hopkins 周期映射。它是满射 étale，纤维为 $\mathrm{GL}_n(E)/K$。特别地，adic 空间 $\PP^{n-1}$ 容许非平凡无限次 étale 覆盖空间！

空间 $\mathcal M_K$ 参数化一维、高 $n$ 的 $\pi$-可除 $\cO_E$-模的形变，外加水平结构。对 $K=\mathrm{GL}_n(\cO_E)\subset\mathrm{GL}_n(E)$，$\mathcal M_K$ 是 $\ZZ$ 个 $(n-1)$ 维开单位圆盘的无交并，从而同构于
\[
\coprod_{\ZZ}\left(\Spa W_E(\ol{\Fq})[[u_1,\dots,u_{n-1}]]\right)_{\breveE}.
\]
特别地，$(n-1)$ 维开单位圆盘可以表示成 $\PP^{n-1}$ 的无限次 étale 覆盖。
\end{example}

\begin{remark}
这些例子相当特殊。例如，Lubin–Tate 情形本质上是周期映射 $\pi_K$ 满射的唯一情形。此外，它们是周期映射的像能被显式理解的少数实例。
\end{remark}

现在谈谈局部 Shimura 簇的构造。我们要一个开子集 $\mathcal{F}\ell^{\mathrm{adm}}_\mu\subset\mathcal{F}\ell_\mu$（「容许轨迹」）以及 $\mathcal{F}\ell^{\mathrm{adm}}_\mu$ 上的 $G(E)$-局部系统 $\mathbb L$。于是可以定义 $\mathcal M_K$ 来参数化 $\mathbb L$ 到 $K$ 的约化；等价地，把 $\mathbb L$ 看作 $G(E)$-主丛
\[
\mathbb L\to\mathcal{F}\ell^{\mathrm{adm}}_\mu,
\]
可以定义 $\mathcal M_K=\mathbb L/K\to\mathcal{F}\ell^{\mathrm{adm}}_\mu$。这将自动是 étale 的，因为 $G(E)/K$ 离散。由于 $\mathcal{F}\ell^{\mathrm{adm}}_\mu$ 在旗簇中是开的，它自然具有光滑刚性解析簇的结构，所以 $\mathcal M_K$ 是光滑刚性解析簇。

所以构造的内容在于找到正确的开子集与正确的 $G(E)$-局部系统。回忆我们已定义 $\Gr_G:=\Gr^{B^+_{\mathrm{dR}}}_G$。在 $\Gr_G$ 内有 Schubert 簇 $\Gr_{G,\le\mu}$；由于 $\mu$ 是极小的，它同构于 $\mathcal{F}\ell^{\diamond}_\mu$。（若 $\mu$ 不是极小的，则得到类似的故事，只是把 $\mathcal{F}\ell^{\diamond}_\mu$ 换成 $\Gr_{G,\le\mu}$。）

\srcpage{69}

我们已经通过修正平凡丛构造了单值化映射 $\Gr^{B^+_{\mathrm{dR}}}_G\to\Bun_G$。但现在想考虑由 $b\in B(G)$ 给出的另一个映射，即修正 $\cE_b$。于是得到映射
\[
\mathcal{F}\ell^{\diamond}_\mu\to\Bun_G.
\]

\begin{proposition}[Rapoport 为 {[S18]} 所写的附录]
$\mathcal{F}\ell^{\diamond}_\mu\to\Bun_G$ 的像与 $\Bun^1_G$ 相交，当且仅当 $b\in B(G,\mu)$。\footnote{此处一切均差一个符号。}
\end{proposition}

所以 $B(G,\mu)$ 正是那些等晶体，使得存在类型 $\mu$ 的修正把 $\cE_b$ 送到 $\cE_1$；等价地，使得存在类型 $\mu^{-1}$ 的修正把 $\cE_1$ 送到 $\cE_b$。也就是说，$\cE_b$ 落在通过修正 $\cE_1$ 得到的类似映射 $\mathcal{F}\ell_{\mu^{-1}}\to\Bun_G$ 的像中。

开子叠 $\Bun^1_G\hookrightarrow\Bun_G$ 于是诱导开子簇 $(\mathcal{F}\ell^{\mathrm{adm}}_\mu)\hookrightarrow\mathcal{F}\ell^{\diamond}_\mu$。
\[
\begin{tikzcd}
\mathcal{F}\ell^{\diamond}_\mu \arrow[r] & \Bun_G \\
\mathcal{F}\ell^{\mathrm{adm}}_\mu \arrow[r] \arrow[u, hook] & \Bun^1_G \arrow[u, hook] \arrow[d, "\sim"] \\
& {[\mathrm{pt}/G(E)]}
\end{tikzcd}
\]
这给出了所要的数据：开子簇 $\mathcal{F}\ell^{\mathrm{adm}}_\mu$，以及其上的一个 $G(E)$-主丛。

\begin{corollary}
无限水平的局部 Shimura 簇
\[
\varprojlim_K\mathcal M^{\diamond}_K
\]
参数化类型 $\mu$ 的修正 $\cE_b\cong\cE_1$。更确切地说，对一切 $S\in\Perf/(\Spa\breveE)^{\diamond}$，
\[
\varprojlim_K\mathcal M^{\diamond}_K(S)=\left\{\begin{array}{c}
\text{同构 }\cE_b|_{X_S\setminus S^{\#}}\cong\cE_1|_{X_S\setminus S^{\#}}\text{ 的}\\
\text{类型 }\mu\text{ 修正（在几何点处）}
\end{array}\right\}.
\]
\end{corollary}

\begin{example}
在上述 Lubin–Tate 情形中，Lubin–Tate 塔的反极限
\[
\varprojlim_{K\subset\mathrm{GL}_n(E)}\mathcal M^{\diamond}_K\cong\{\cO^n_{X_S}\hookrightarrow\cO_{X_S}(1/n)\text{ 的余核支撑在 }S^{\#}\}.
\]
注意余核由于次数原因必然是 $S^{\#}$ 上的线丛。

也可以研究 Drinfeld 塔
\[
\varprojlim_{K'\subset D^{\times}_{1/n}}\mathcal M^{\diamond}_K,
\]
结果表明它在 $\infty$ 水平有同样的描述，即作为 $\cO(1/n)$ 与平凡丛之间的修正（但顺序相反）。
\end{example}

\begin{remark}
这一分析推广后，对 $b$ 基本时的所有局部 Shimura 簇给出对偶。它联系 $(G,b,\mu)$ 与 $(G_b,b^{-1},\mu^{-1})$。
\end{remark}

\begin{remark}
对一般的 $\mu$，有 Białynicki–Birula 映射
\[
\Gr_{G,\mu}\to\mathcal{F}\ell^{\diamond}_\mu.
\]

\srcpage{70}

纤维是迭代仿射空间（但它本身不是仿射空间，因为存在非分裂扩张 $(\AA^1)^{\diamond}\to X\to(\AA^1)^{\diamond}$，其中 $X$ 不由刚性解析簇表示）。

设 $S$ 是刚性解析簇，为简单起见设其光滑，并考虑映射 $S^{\diamond}\to\mathcal{F}\ell^{\diamond}_\mu$。可以问它到 $\Gr_{G,\mu}$ 的提升。
\[
\begin{tikzcd}
& \Gr_{G,\mu} \arrow[d] \\
S^{\diamond} \arrow[ur, dashed] \arrow[r] & \mathcal{F}\ell^{\diamond}_\mu
\end{tikzcd}
\]
你也许会猜有很多提升，但事实上并非如此。
\end{remark}

\begin{theorem}[{[S14]}]
映射 $\Gr_{G,\mu}(S^{\diamond})\hookrightarrow\mathcal{F}\ell^{\diamond}_\mu(S^{\diamond})=\mathcal{F}\ell_\mu(S)$ 是单射，其像由满足 Griffiths 横截性（transversality）的映射组成。\footnote{这在一定程度上要归功于与 Kedlaya 的讨论；Fontaine 与 Faltings 可能早已知道。}
\end{theorem}

这当 $S$ 是一个点时已经很有趣。

\subsection{$\Bun_G$ 的一般点}

\subsubsection{半稳定点}

\begin{theorem}
半稳定轨迹 $\Bun^{\mathrm{ss}}_G\subset\Bun_G$ 是开的，且
\[
\Bun^{\mathrm{ss}}_G=\coprod_{b\in B(G)_{\mathrm{basic}}}\underbrace{[\mathrm{pt}/G_b(E)]}_{\Bun^b_G},
\]
其中 $\Bun^b_G$ 是几何点上 $\cE\cong\cE_b$ 的轨迹。
\end{theorem}

\begin{proof}
关于开性的陈述由 $\nu$ 的半连续性推出。$\Bun^b_G$ 中的分解由 $\kappa\colon|\Bun^b_G|\to B(G)$ 的局部常值性以及 $\kappa$ 限制到基本轨迹上是同构 $B(G)_{\mathrm{basic}}\xrightarrow{\sim}\pi_1(G)_{\Gamma}$ 推出。所以只需证明 $\Bun^b_G\cong[\mathrm{pt}/G_b(E)]$。

但注意 $\cE_b$ 是 $X_S$ 上的 $G$-主丛，且对基本的 $b$ 有 $\Aut_{X_S}(\cE_b)=G_b\times_E X_S$。（一般地，自同构群会是一个内形式，因为容易检验它局部同构于 $G$。）事实上，$G_b$ 以显然的方式通过自同构作用在 $\cE_b$ 上；由于它是 $G$ 的内形式，它必须填满所有自同构。特别地，得到等价
\[
\{X_S\text{ 上的 }G\text{-主丛}\}\cong\{X_S\text{ 上的 }G_b\text{-主丛}\},
\]
把 $\cE$ 送到 $\underbrace{\Isom(\cE,\cE_b)}_{\Aut(\cE_b)\text{-主丛}}$。
因此基本的 $b$ 诱导同构 $\Bun_G\cong\Bun_{G_b}$，把 $\Bun^b_G$ 同构地送到 $\Bun^1_{G_b}$。
\[
\begin{tikzcd}
\Bun_G \arrow[r, "\sim"] \arrow[d] & \Bun_{G_b} \arrow[d] \\
\Bun^b_G \arrow[r, "\sim"] & \Bun^1_{G_b}=[\mathrm{pt}/G_b(E)]
\end{tikzcd}
\]
\end{proof}

\srcpage{71}

\subsection{非半稳定的 $b$}

\begin{theorem}
$\Bun^b_G\subset\Bun_G$ 局部闭，且
\[
\Bun^b_G\cong[\mathrm{pt}/\mathcal G_b],
\]
其中 $\mathcal G_b$ 是群 $v$-层，它落在短正合列
\[
1\to\mathcal G^0_b\to\mathcal G_b\to G_b(E)\to 1,
\]
且 $\mathcal G^0_b$ 是正 Banach–Colmez 空间的扩张，维数为 $\langle 2\rho,\nu(b)\rangle$。\yizhu{原文此处用花体 $\mathcal G_b$，与前面定理中的 $G_b$ 不同。}{A（原文 p.71）}
\end{theorem}
